% 06-conclusion-limitations.tex — truth: 本文件（2026-09-10 交割：R22 后全部正文修订直改 tex，原装配源 md 已冻结为历史设计稿，见其头部治理注记）
% 纪律: 结论段逐字搬运; Limitations 十二条独立成段 (L1-L12; 终审 2026-09-07 勘正条数); rebuttal 预案不入正文(投稿期备用)。
\section{Conclusion and limitations}\label{sec:conclusion}

\subsection{Conclusion}

We presented PSD, a physics--semantics decoupled framework for animal behavior recognition under evolving evaluation criteria. PSD separates how skeletons move from what behaviors are called: a label-free physics layer provides self-supervised dynamics representations and unsupervised behavior proposals, while a lightweight semantic layer grows rule-engine seeds toward full taxonomy coverage through anchor-guided clustering, iterated pseudo-labeling, and warm-started small-budget initialization. Experiments across the two exercised tiers and four public benchmarks spanning two kingdoms (NTU60, NTU120, UCF101, PanAf500) showed that pretraining alone yields discriminable dynamics representations ($2.51\times$ random baseline), that the semantic layer attains \textbf{82.0\% top-1 on 22 classes from a 20-clip budget via warm-started initialization} (synthetic-offset tier; $\pm$4.3 across three seeds), and that taxonomy transitions are absorbed at \textbf{$\geq 3\times$ lower wall-clock retraining cost} than full-pipeline retraining at matched accuracy (conservative bound backed by a measured $6.07\times$ on the synthetic tier; accuracy differences inside the pre-registered $<2.3$\,pp noise band ($-0.91$\,pp full / $+2.27$\,pp small)). An implementation-equivalence verification on NTU60 (three-stream fusion 77.97\% against a pre-registered acceptance line of 77.18\%) shows that the re-implemented pipeline meets its pre-registered equivalence criterion against the official reference it adapts. Under a corrected pseudo-label protocol the end-to-end pipeline retains 90.7\% of full-budget accuracy at 10\% of the labels on the human benchmark (NTU120 88.9\% and PanAf500 88.7\% hold the PARTIAL band, with UCF101 at 66.6\% marking the collapse boundary; a linear head on frozen features alone retains 88.7\% on NTU60, and the pseudo-label iteration adds $+1.48$\,pp, clearing the pre-registered 90\% line), while at a 13\% budget on the canine tier the iterated pseudo-labels degrade and accuracy stays near chance---a tier-dependent boundary that bounds the low-resource claim to the regimes measured. Beyond animal behavior, the decoupling pattern---frozen physics, revisable semantics---applies to recognition tasks whose evaluation criteria evolve with operational practice---we verify the pattern on skeleton data in two domains and state it as a design claim elsewhere.

\subsection{Limitations}

Thirteen limitations bound our claims.

\paragraph{L1: Proxy-caliber evidence for representation quality.}
Our pretraining evidence uses subject identity as a $k$-NN probe on InterPet4D, which measures representation discriminability rather than behavior accuracy; behavior-level evidence rests on the downstream experiments of Sections~\ref{sec:experiments}--\ref{sec:ablation} (segmentation IoU 0.458$\pm$0.049 under a seeds pseudo-ground-truth protocol, pseudo-label pool precision 0.691$\pm$0.013, the four-class public-real fine-tuning result below, and the cross-domain end-to-end retention of E9), each with its caliber disclosed inline.

\paragraph{L2: Single species family.}
The public-real tier's real-data validation covers canids only; all other real-data tiers are validated on their own domains (the PanAf500 great-ape benchmark is validated in-domain, E9d), and no cross-family pretext transfer (canid$\to$felid/equid/primate) has been tested---the NTU human-domain verification is not animal cross-family evidence and is left to future work. Sex- and gender-based analyses were not performed (see the Ethics statement); no dataset in this study supports them.

\paragraph{L3: Survey boundary on firstness.}
Our novelty claim rests on repository-scale search with zero occupancy; an arXiv/CrossRef re-verification at submission time surfaced point-supervised skeleton segmentation as the nearest neighbor (differentiated in Section~\ref{sec:related}), and we cannot exclude unpublished or non-indexed concurrent work.

\paragraph{L4: Scale of compute.}
All experiments ran on a single consumer GPU (RTX 5060 Laptop 8\,GB); scaling behavior to multi-GPU pretraining and web-scale unlabeled pools is unverified.

\paragraph{L5: No advantage of uncertainty sampling---reproduced across three independent runs.}
In our short-budget active-learning studies (synthetic tier), entropy sampling did not outperform random selection under either weak or strong scorers: with cold-start scorers at the short schedule, random led by 7.9\,pp at budget 100 (2 of 3 seeds) and 7.1\,pp at 200 (3 of 3 seeds); a full-budget cold-start rerun (120 vs.\ 50 training epochs) reproduces the direction (random ahead by 12.3\,pp at budget 100 and 7.8\,pp at 200); with warm-started in-domain fine-tuned scorers, random again led by 4.2--5.0\,pp across all three budget points. The synthetic-trained checkpoint additionally scored the public-real candidate pool with fully saturated softmax (mean top1$-$top2 logit margin 100.9 vs.\ $\approx$10.8 within the synthetic validation domain), degrading uncertainty signals to numerical zero. We therefore report uncertainty sampling strictly as an exploratory negative finding and make no sampling-efficiency claims.

\paragraph{L6: Structural constraints of the public-real tier.}
Under the pre-registered partial-class protocol, only 4 of 12 target classes pass the canine sample-size gate under the relaxed rule (3 under the strict rule; down/stand/scale have zero coverage), and in the released pose-estimation subset we measure $\approx$4.6 frames per video (28{,}197 images across 6{,}160 video directories; the Animal Kingdom paper reports 33{,}099 pose frames, a counting difference between the official figure and our local extraction, disclosed here), with no canine-overlapping video reaching the $T{=}30$ temporal length required by skeleton pipelines; a self-extraction pipeline (YOLO11-pose fine-tuned on dog-pose, then skeleton extraction and ST-GCN head fine-tuning on the 4-class subset) was adopted as mitigation, and its first-round outcome is quantified in L7.

\paragraph{L7: Severe class imbalance dominates the public-real result.}
The self-extraction round-one model reaches 44.9\% overall validation accuracy on the four-class subset (seed 42; three-seed mean 41.5\%, sample std $\pm$5.9\,pp; $1.80\times$ the 25\% random baseline) under heavily skewed support (full-list: watch 72 / track 46 / stay 27 / jump 27; train split: 57/29/19/18); per-class accuracies collapse to watch 100\% / track 23.5\% / jump 0\% / stay 0\%---majority-class recall accounts for essentially all of the aggregate score. We therefore present this number with its per-class breakdown attached, restrict this aggregate's interpretation to the four-class partial caliber, and treat aggregate-only reporting as prohibited for this tier.

\paragraph{L8: Real-domain skeletons carry 20/24 valid supervision channels.}
The dog-pose annotations never label four of the 24 graph joints (both eyes, withers, throat---the last four slots are absent from the source data), so every real-domain keypoint product in this paper is topology-isomorphic rather than fully supervised: the four dead channels are hard-masked at assembly and NaN-guarded downstream. Any real-domain skeleton-based number therefore rests on 20-channel effective supervision, and rules that depend on the withers landmark are structurally degraded in the pixel domain.

\paragraph{L9: Boundary conditions of cross-domain unsupervised enhancement---pre-registered negative at this resource level.}
We tested whether unlabeled cross-domain geometry transfers into the working-dog 2D pose domain of the public-real tier: a pre-registered round augmented the frozen backbone with AdaBN-style statistics from 1{,}069 clips (570 same-pipeline pixel-domain extraction, 499 heterogeneously annotated APTv2 domain) under a kinematic plausibility gate calibrated on 51 motion-capture sequences, then repeated the identical four-class head re-training protocol. The primary endpoint failed: 40.82\% vs.\ 44.90\% ($-4.08$\,pp against a bit-identical same-day baseline rerun; paired three-seed deltas $-4.08$/$+6.12$/$-8.16$\,pp, mean $-2.04$\,pp inside the $\pm$5.9\,pp seed-noise band). Because the adaptation demonstrably penetrated the backbone (85/85 BatchNorm statistics moved; the track profile swung 23.5\%$\to$76.5\% in the strongest single-source arm), the bottleneck is conversion, not supply: absent a label-alignment force between unsupervised second-order statistics and small-sample head re-training, no adapted arm yields a legitimate aggregate gain---the only same-seed aggregate rise ($+4.08$\,pp, APTv2-only arm) is a class-level swap (track tripled while watch fell 100\%$\to$26.7\%) reported as a secondary signal only. Motion-capture kinematics transfer only through a one-sided filter: the treadmill-only reference set contains no static behavior, so its low-speed band encodes reference bias rather than physics. We therefore claim no accuracy benefit from unlabeled cross-domain enhancement at this caliber and identify label-in-the-loop alignment, not more unlabeled volume, as the missing fuel for any future flywheel claim---a diagnosis that the corrected end-to-end arms of E7/E7b quantify on the canine tier (pseudo-label precision $\approx$0.11 at a 13\% budget; near-chance end-to-end accuracy).

\paragraph{L10: The pretraining-benefit crossover point is bracketed, not localized.}
A low-resource gradient ablation (synthetic tier; scratch vs.\ pretraining-initialized, two arms $\times$ three seeds, strictly matched budgets) shows the downstream fine-tuning benefit of self-supervised pretraining decaying with annotation budget to statistical equivalence (with a disclosed $-2.3$\,pp reversal at 352 clips): $+21.2$\,pp at 88 training clips, $+9.9$\,pp at 176, $-2.3$\,pp at 352, and $+0.15$\,pp (n.s.) at 1760. The zero crossing lies between 10 and 20 clips per class---linear interpolation places it near 15--20---but no finer tiers were run, so its exact location remains unlocalized. Two caliber constraints apply: the gradient tiers ran on CPU while the saturation tier ran on GPU, so cross-tier absolute accuracies are not poolable (trend claims rest on within-tier paired deltas only); and the lowest tier validates on only 22 samples (one sample = 4.55\,pp), so effect sizes there are directional evidence rather than precise estimates. The curve is presented whole; no claim extends outside the measured tiers.

\paragraph{L11: E9's retention cannot separate filtering from pool size.}
At NTU scale the confidence gate accepts $\approx$98\% of the unlabeled pool, so the 90.7\% retention characterizes the pipeline's budget behavior under pseudo-label restoration and not the contribution of filtering per se; the 10\% labeled subset is a single stratified draw, so the reported spread covers self-training seeds but not subset resampling; and the loop's precision-drop stopping rule is disabled on this tier (no rule-derived consensus reference exists), with pool precision against the official labels ($\approx$0.69) reported only as a post-hoc diagnostic outside the control path. The external-SSL comparison is likewise harness-limited: the Mean-Teacher baseline runs on the same frozen features with a confidence threshold ($\tau{=}0.95$) whose selectivity is the opposite of the pipeline's near-full pool acceptance, so the $+20.7$\,pp gap measures the value of near-full pool restoration over selective pseudo-labeling, not the superiority of pseudo-label restoration over SSL in general---an end-to-end SSL comparison (video-level, learned representations) remains outside this study's scope. A post-hoc $2\times2$ decomposition (PSD-POOLDECOMP-PREREG-001, registered after the series; four arms differing only in pool composition---gate on/off $\times$ full pool/random equal-size subset, ten seeds each) measured both factors directly: the filtering main effect is $-0.02$\,pp and the pool-size main effect is $0.0$\,pp (all four arms within $67.53$--$67.58\%$), so at this scale the pipeline's $+1.48$\,pp gain is attributable to pseudo-label self-training over a near-full pool, with neither gate quality nor pool size contributing measurably---the mechanism reading asserted above is confirmed experimentally, and the gate is inert at this operating point (it is load-bearing only in the canine tier's selective regime, where it filters on calibrated margins).

\paragraph{L12: The $\geq 3\times$ transition-cost bound.}
The bound is measured on the synthetic tier (E6, $6.07\times$; seeds 42--44) and replicated on the human benchmark at backbone scale (PSD-NTU-TRANS-002, frozen before the run, seeds 42--44): at the full budget the median coupled/decoupled wall-clock ratio is 15.8x (9897\,s vs 626\,s; a final-segment-only bound of 6.2x still clears the 3x line) with a coupled-over-decoupled accuracy margin of 4.5\,pp (outside the $\pm$2.3\,pp band), and at the 10\% budget the ratio reaches 32.5x with a +2.6\,pp decoupled advantage (PARTIAL at both budgets). We did not measure a real-data taxonomy transition on working-dog data because no public canine dataset provides two aligned taxonomies over the same unlabeled stream; the bound's magnitude on real working-dog data is an open question the pre-registered K9 pilot addresses. The MLP-scale TRANS-001 failure is therefore re-read as scale-bound (no cost to save at 35\,s), while the backbone-scale replication shows the cost saving is real but the matched-accuracy condition holds only approximately. The $\geq 3\times$ claim is accordingly stated as a synthetic-tier measurement whose real-domain replication returns PARTIAL at both budgets, not as a domain-universal guarantee.

\paragraph{L13: No head-to-head comparison against any published method.}
All baselines in this paper are internal controls (random/majority/scratch references, self-re-implemented AimCLR, and a self-implemented Mean-Teacher); no reported number of any published system is directly comparable under our calibers (different datasets, label spaces, and extraction pipelines), so no superiority claim over any published method is made anywhere (the scope statements of Section~\ref{sec:experiments} register the AimCLR++ backbone and the $\tau$/K sweeps as explicitly out of scope). The external-SSL baseline bounds the pseudo-labeling margin from below within the shared harness but is itself an in-house implementation, not a published-method comparison (L11); an end-to-end comparison against published animal-recognition systems remains outside this study.
